% ======================================================================
% 2.1 Поворот декартовых координат. Ответственный: D — Филиппов Андрей.
%
% Что сдать (main.tex: условие 436–456, порядок 457–487, пример 628–668;
% требования к отчёту 1862–1870):
%   1. Изображения исходной области (x,y,z) и новой (u,v,w), оси и границы.
%   2. Матрица частных производных обратной замены, её определитель, проверка |J|=1.
%   3. Подстановка x,y,z в функцию, упрощение, повторный интеграл по u,v,w
%      с явными пределами и множителем якобиана.
%   4. Аналитическое вычисление последовательным интегрированием.
%   5. Метод средних прямоугольников при n=8,16,32,64; шаги h1=a/n, h2=b/n, h3=c/n.
%   6. Проверка объёма: интеграл от 1 (с якобианом) сравнить с V=abc.
%   7. Таблица значений и ошибок, график ошибки от h=1/n, вывод о сходимости.
%
% Параметры варианта 3 (приложение А): a=2, b=3, c=1, alpha=pi/3.
% Код: solution.py (блок D), запуск python experiments.py 2.1.
% Рисунки: figures/2-1-*.pdf.
% ======================================================================

\subsection{Поворот декартовых координат}

\taskhead{Исходные данные и аналитическое решение}
Для варианта 3 заданы $a=2$, $b=3$, $c=1$, $\alpha=\pi/3$. Рассматривается интеграл
\[
I_{\mathrm{п}}=\iiint_{\Omega_{\mathrm{п}}}(x^2+y^2+z^2)\,dx\,dy\,dz,
\]
\[
\Omega_{\mathrm{п}}=\{0\leq x\cos\alpha+y\sin\alpha\leq a,\quad
0\leq -x\sin\alpha+y\cos\alpha\leq b,\quad 0\leq z\leq c\}.
\]
\todo{D: обосновать выбор координат; выписать поворот и обратную замену с параметрами варианта; составить матрицу частных производных обратной замены и вычислить определитель (проверить $|J|=1$); подставить $x,y,z$ в функцию и упростить; записать повторный интеграл по $u,v,w$; вычислить аналитически (пример: main.tex, строки 628–668).}

\taskhead{Программный результат и сравнение}
\todo{D: листинг вызова своей функции из solution.py; указать шаги $h_1=a/n$, $h_2=b/n$, $h_3=c/n$ для каждой сетки.}
Аналитическое значение $I_{\mathrm{п}}$: \todo{D: подставить}.
На графике используется нормированный шаг $h=1/n$ (main.tex, строки 1989–1990, 2078–2082).
\begin{table}[htbp]
\centering\small
\caption{Сравнение численного интеграла с аналитическим значением $I_{\mathrm{п}}$}
\label{tab:2-1}
\begin{tabular}{|r|r|r|r|r|}\hline
$n$ & $h$ & $I_n$ & $\Delta_n$ & $\delta_n$\\\hline
8 & — & — & — & —\\
16 & — & — & — & —\\
32 & — & — & — & —\\
64 & — & — & — & —\\
\hline\end{tabular}
\end{table}

\taskhead{Проверка объёма}
\todo{D: записать интеграл от 1 по исходной области, преобразовать с учётом якобиана, вычислить численно на всех четырёх сетках и сравнить с $V=abc$.}

\taskhead{Графическая проверка}
\taskfigure[0.55]{figures/2-1-region-old.pdf}{Исходная область в координатах $(x,y,z)$ (равный масштаб осей)}{fig:2-1-region-old}
\taskfigure[0.55]{figures/2-1-region-new.pdf}{Новая область в координатах $(u,v,w)$ (равный масштаб осей)}{fig:2-1-region-new}
\taskfigure[0.75]{figures/2-1-error.pdf}{Абсолютная ошибка от нормированного шага $h=1/n$ (логарифмические оси)}{fig:2-1-error}
\todo{D: пояснить график.}

\taskhead{Вывод}
\todo{D: вывод о сходимости; объяснить поведение ошибки при удвоении $n$.}
